\cotefile{156-9}
\scanpage{66}
\mspage{62}
% source lines 1--6
que si de plus $\alpha$ est \unsure{surjectif}, i.e. \unsure{bijectif}!

\unsure{Donc} il faut une condition supplémentaire pour pouvoir affirmer la commutation de $\alpha_{\unsure{*}}$ à $C$, ou mieux, la formule $(**)$. On va supposer donné une correspondance en sens inverse
\diagram{\begin{tikzcd}[ampersand replacement=\&]
E' \arrow[r,"\beta"] \& \unsure{\mathcal{P}|\mathscr{P}}(E)
\end{tikzcd}}
% source lines 7--10
\unsure{i.e.} une application
\diagram{\begin{tikzcd}[ampersand replacement=\&]
\mathcal{P}(E') \arrow[r,no head,"\bar\beta"] \& \mathcal{P}(E)
\end{tikzcd}}
commutant aux réunions. On va supposer que
\[
\forall x'\in E',\quad \bar\alpha\bigl(\beta(x')\bigr)\subset\operatorname{supp}^{\circ}_{E'}(\{x'\}). \tag*{$(*)$}
\]
% source lines 11--16
Si $B\subset E$, \deleted{alors} et $x'\in E'$, alors
\[
x'\mid\circ\mid\bar\alpha(B)\Longleftrightarrow\bar\alpha\bigl(\beta(x')\bigr)\mid\circ\mid\bar\alpha(B)
\]
% margin note written slanted
\margin{\uline{NB} \unsure{si l'hypoth.} \unsure{$(*)$} \unsure{est une hyp.} \unsure{de fidélité}}
et, si $\bar\alpha$ a une propriété de \uline{fidélité} pour $\mid\circ\mid$, on en conclut
\[
\beta(x')\mid\circ\mid B
\]
% source lines 17--21
Appliquons ceci à la situation $(**)$ (p. \unsure{précéd.}) à $(\unsure{B},x')$ et à $(C,y')$, on trouve
\[
\beta(x')\mid\circ\mid B;\quad
\underbrace{\beta(y')\mid\circ\mid C=\operatorname{cosupp}^{\circ}(B)}_{\text{i.e.}\quad\beta(y')\subset\operatorname{supp}^{\circ}(B)}
\]
d'où
\[
\beta(x')\mid\circ\mid\beta(y'),\qquad\text{d'où}
\]
% archive sheet number: 65
